% Standalone manuscript input; not included in the main paper.
\section{Two local obstructions to lifting the characteristic-three bank}
\label{sec:charthree-local}
The results here concern specified collision points of incidence schemes.
They do not bound lists at arbitrary points or exclude other constructions
with eight candidates.

Let $k=\mathbb F_{3^{18}}$, choose a primitive seventh root $\zeta$, and put
\[
 \eta=\zeta+\zeta^2+\zeta^4,\quad
 \bar\eta=-1-\eta,\quad \alpha=\bar\eta/\eta,\qquad
 P(X)=\eta X^3+X^2+X+\bar\eta.
\]
Define $P_i(X)=\zeta^{5i}P(\zeta^{-i}X)$ for $0\le i<7$,
and $P_7=0$. On the fourteen points
$\zeta^t,\alpha\zeta^t$ ($0\le t<7$), take received values
$\zeta^{5t},0$, respectively. Each point has four agreeing candidates;
each candidate agrees at seven points and each pair at three.
Set $\phi(U)=U^3-U$. Its three preimages of every base point lie in $k$
and are distinct, giving forty-two core coordinates.

For $d=1,2$, the core candidates and received word are
$F_i(U)=U^dP_i(\phi(U))$ and $U^dw(\phi(U))$.
For $d=1$, append two coordinates at zero, assigning complementary
four-subsets of candidates. We consider the five partitions whose first
part is $\{7\}\cup T$, where the bit masks of $T\subseteq\{0,\ldots,6\}$
are $7,11,13,19,21$. For $d=2$, append six coordinates at zero and label
candidates by the binary expansions of $0,\ldots,7$. The six prescribed
four-subsets are the coordinate faces
$\{i:\operatorname{bit}_j(i)=b\}$, $0\le j<3$, $b\in\{0,1\}$.
All appended received values are zero.

\paragraph{Local conclusions.}
For each specified simple-factor point, the completed incidence ring over
$W(k)$ is $k[[t_1,\ldots,t_{20}]]$. In particular, there is no
mixed-characteristic formal branch, of any ramification degree, and the two
appended coordinates remain equal in every equal-characteristic formal
deformation. For the specified double-factor point, no DVR-valued branch
specializing to it has all six appended coordinates distinct. This second
statement holds in mixed characteristic and in characteristic three.
The parameters are respectively $(n,D,A)=(44,10,22)$ and $(48,11,24)$;
their apparent quarter-rate, half-agreement configurations therefore do
not lift to distinct-node examples through these points.

\paragraph{Finite certificates.}
The simple-factor Jacobians have rank $156$ in $176$ variables. A
$20$-dimensional kernel, nonzero maximal minors, and nonzero pairings of a
left-null vector with the characteristic defect divided by three were
verified independently. The double-factor Jacobian has rank $164$ in
$192$ variables, with a $28$-dimensional kernel. Its projected quadratic
correction has rank $11$; a verified functional annihilates its entire
span but not the characteristic defect. The correction for an incidence
at $T$ is
\[
 F_i^{[2]}(T)\xi^2+R_i'(T)\xi,
\]
where $R_i$ is the tangent polynomial and brackets denote a Hasse
derivative. Arithmetic was replayed in
$\mathbb F_3[t]/\Phi_7(t^3-t)$, independently of the original tower-field
implementation. The certificates and replay scripts are in the
\texttt{char3\_deformation} and \texttt{char3\_double\_deformation}
research directories. These are exact finite-field identities, not
numerical rank estimates.

\paragraph{The simple-factor formal argument.}
Writing $P_i^{\rm hom}$ for cubic homogenization, consider
\[
 C(U)+L(U)P_i^{\rm hom}(A(U),B(U)),\qquad
 \deg C\le10,\quad \deg L\le1,\quad \deg A,\deg B\le3,
 \quad B(0)=1.
\]
There are twenty parameters. The core coordinates are the formal roots
of $A-aB$, and both appended coordinates are the root of $L$.
This characteristic-three family has injective differential: vanishing
core motions force $\delta A-A\delta B$ to vanish at forty-two points,
hence identically; normalization gives zero rational-map variation.
The zero candidate then forces $\delta C=0$, and the remaining candidates
force $\delta L=0$.
Solve the $156$ pivot equations over $W(k)$ by the formal implicit-function
theorem. The twenty remaining parameters are formally parametrized by this
family modulo three, so every residual equation is divisible by three.
The verified obstruction modulo nine makes one residual equation divided
by three a unit. Their ideal is therefore exactly $(3)$, proving the
claimed completed-ring description.

\paragraph{Constant rank for the double-factor family.}
The corresponding characteristic-three family is
\[
 C(U)+l(U-b)^2P_i^{\rm hom}(A(U),B(U)),\qquad
 \deg C\le11,\quad \deg A,\deg B\le3,\quad B(0)=1.
\]
It has twenty-one parameters and injective differential. Its Jacobian
rank is constantly $164$ on the open set where the degree-three map
$\pi=A/B$ is unramified over the base nodes and $b$ avoids their preimages.
Here is a cover-independent justification. Tangent polynomials satisfy
$R_i(b)=c$ for a common scalar, hence
$R_i=c+(U-b)S_i$, $\deg S_i\le10$. At each core point the conditions on
$(S_i)$ are precisely the fixed base conditions of lying in the span of
the received-value direction and the base derivative direction. Let
$\mathcal E\subseteq\mathcal O^8$ be the elementary modification imposing
these two linear constraints at each of the fourteen base nodes, and put
$h_j=h^0(\mathcal E(j))$. The six appended node velocities are free, and
all core node and word velocities are uniquely determined. Thus nullity is
\[
 7+h^0\bigl(\mathcal O(10)\otimes\pi^*\mathcal E\bigr)
 =7+2h_3+h_2.
\]
Indeed $\pi_*\mathcal O(1)=\mathcal O^2\oplus\mathcal O(-1)$:
its rank, degree, global sections, and vanishing after twisting by
$\mathcal O(-1)$ force this splitting. Projection formula gives
$\pi_*\mathcal O(10)=\mathcal O(3)^2\oplus\mathcal O(2)$.
The nullity is consequently independent of the cover and equals the
certified value $28$.

\paragraph{The normal cone and collision argument.}
After eliminating the $164$ pivot variables and straightening the
$21$-parameter family, use normal coordinates $z_1,\ldots,z_7$.
Constant rank gives residual equations of the form
\[
 g(u,z)=3h(u)+3L(u)z+Q_u(z)+O(z^3).
\]
The normal tangent coordinates can be taken as
$(\omega_0,\ldots,\omega_5,\tau)$, where
$\omega_f=\xi_f+\alpha_1/2$ and the polynomial slopes are
$\beta+\alpha_1c_i+\tau d_i$. Here $c_i=[U^2]F_i$ are pairwise distinct.
Eleven explicit row combinations of the original projected tensor,
independently replayed coefficient by coefficient, give
\[
 \omega_i^2-\omega_5^2=0,\qquad
 \tau(\omega_i-\omega_5)=0\quad(0\le i<5),\qquad \tau^2=0.
\]
These seven coordinates vanish on the family tangent space and form its
full normal quotient.

For a mixed-characteristic branch with uniformizer $\varpi$, write $e=v(3)$ and
$s=\min_jv(z_j)$. The case $z=0$ is excluded by $h(0)\ne0$.
If $e<2s$ the characteristic forcing cannot cancel; if $e=2s$ the verified
separating functional prevents cancellation. Thus $e>2s$ and the leading
normal vector lies on the displayed cone. It is nonzero, so
$\tau=0$ and $\omega_f\in\{a,-a\}$ with $a\ne0$.
Centering at $b-\varpi^s\alpha_1/2$ eliminates the order-$s$ linear term in
pair differences. Incidence at a shared mark forces their constant term
to have order at least $2s$. Consequently, for adjacent cube vertices,
\[
 \overline{\varpi^{-2s}(F_i-F_j)(\mathrm{center}+\varpi^sZ)}
       =(c_i-c_j)(Z^2-a^2).
\]
The unknown second-order constants are determined here by evaluation at
the shared marked velocity; they are not assumed to vanish.

Two faces from different axes share an adjacent candidate pair. If they
have the same sign velocity, simple-root uniqueness for this reduced
quadratic forces their actual coordinates to coincide. Hence distinct
marks in either sign cluster can belong to only one axis, at most two
faces. Two clusters cannot contain six distinct marks. In equal
characteristic, $z=0$ gives coincident marks on the family; otherwise the
same homogeneous-cone and simple-root argument applies directly.
